
The child-to-parent measures are
$du/[u(1-u)]$ and $du\,dv/[uv(1-u-v)]$.
They are positive.  To specify the one-step switched cost, define
\begin{equation}\label{sv:one-step-mass}
 I(s,s_0,\alpha)=
 \int_{1/s_0<u_1<u_2<u_3<1/s}
 \omega\!\left(\frac{\alpha-u_1-u_2-u_3}{u_2}\right)
                    \frac{du_1\,du_2\,du_3}{u_1u_2^2u_3}.
\end{equation}
The one-step inequality gives the candidate
\begin{align}\label{sv:one-step-operator}
 \mathcal W_1[H,h](s;s_0,\alpha)
 ={}&A_+(s)-A_+(s_0)+H(s_0)-\frac K8 I(s,s_0,\alpha)\notag\\
 &+\frac12\int_{1/s_0}^{1/s}
 \frac{A_-(s_0(1-u))+h(s_0(1-u))}{u(1-u)}\,du,
 \qquad 2\le s<s_0\le5.
\end{align}
Here $A_+(s)=sF(s)/(2e^\gamma)$ and
$A_-(s)=sf(s)/(2e^\gamma)$, with the zero extension below
$2$.  The physical singleton child cutoff is $N^{L/s_0}$ and its
argument is $s_0(1-u)$, not the argument at a different total level.
The positive switched family is
$Bq_1q_2q_3n\le N$ with
$N^{L/s_0}\le q_1<q_2<q_3<N^{L/s}$ and $P^-(n)\ge q_2$.
Its unsifted mass, in the parent normalization, is $I/4$; the
factor $1/2$ in the pointwise weighted inequality gives the cost
$KI/8$ in \eqref{sv:one-step-operator}.  All its prime factors
have fixed positive exponents.  Its unit quotient has a fixed power
